% !PRINTED 308
\weberpage{308}
%<<CONT
\weq{2}{
\begin{aligned}
\tau &= \begin{pmatrix} 1, & 0, & 1 \\ 1, & 0, & 0 \\ 0, & 1, & 0 \end{pmatrix}, &
\tau^2 &= \begin{pmatrix} 1, & 1, & 1 \\ 1, & 0, & 1 \\ 1, & 0, & 0 \end{pmatrix}, &
\tau^3 &= \begin{pmatrix} 0, & 1, & 1 \\ 1, & 1, & 1 \\ 1, & 0, & 1 \end{pmatrix}. \\
\tau^4 &= \begin{pmatrix} 1, & 1, & 0 \\ 0, & 1, & 1 \\ 1, & 1, & 1 \end{pmatrix}, &
\tau^5 &= \begin{pmatrix} 0, & 0, & 1 \\ 1, & 1, & 0 \\ 0, & 1, & 1 \end{pmatrix}, &
\tau^6 &= \begin{pmatrix} 0, & 1, & 0 \\ 0, & 0, & 1 \\ 1, & 1, & 0 \end{pmatrix}
\end{aligned}
}

Hierauf sucht man unter den Elementen 2\textsuperscript{ten} Grades, deren
es 21 giebt, ein passendes f\"ur $\omega$ aus; es eignet sich z.\,B. dieses

\weq{3}{\omega = \begin{pmatrix} 1, & 0, & 0 \\ 0, & 1, & 0 \\ 0, & 1, & 1 \end{pmatrix};}

und dann sind nach den Formeln \S.\,72, (17) die Elemente $\Theta$
und $\chi$ bestimmt:

\weq{4}{\Theta = \begin{pmatrix} 1, & 1, & 1 \\ 0, & 1, & 0 \\ 0, & 1, & 1 \end{pmatrix}, \quad \Theta^2 = \begin{pmatrix} 1, & 1, & 0 \\ 0, & 1, & 0 \\ 0, & 0, & 1 \end{pmatrix}, \quad \Theta^3 = \begin{pmatrix} 1, & 0, & 1 \\ 0, & 1, & 0 \\ 0, & 1, & 1 \end{pmatrix}}

\weq{5}{\chi = \begin{pmatrix} 1, & 0, & 0 \\ 0, & 0, & 1 \\ 0, & 1, & 1 \end{pmatrix}, \quad \chi^2 = \begin{pmatrix} 1, & 0, & 0 \\ 0, & 1, & 1 \\ 0, & 1, & 0 \end{pmatrix}.}

Hierauf ist es Sache einer einfachen Rechnung, die Formeln
des Theorems I., \S.\,72 zu best\"atigen, wodurch die Uebereinstimmung
 der Gruppen $S$ und $L_7$ nachgewiesen ist.

Es folgt daraus, dass die Gruppe $S$ ebenso wie $L_7$ \gesperrt{einfach}
 ist.

\begin{center}
\S.\,80.

\gesperrt{Resolventen der Gleichung 8\textsuperscript{ten} Grades.}
\end{center}

Um nun die Anwendung auf die Theorie der Gleichungen
8\textsuperscript{ten} Grades zu machen, bezeichnen wir wie fr\"uher mit $Q$ die
cyklische Gruppe 8\textsuperscript{ten} Grades: % [sic] the displayed translation group is elementary abelian of order 8, not cyclic.

\[
\begin{pmatrix} z_1, & z_2, & z_3 \\ z_1 + h_1, & z_2 + h_2, & z_3 + h_3 \end{pmatrix} \text{ (mod 2)},
\]

oder k\"urzer

\[
(z_i,\ z_i + h_i),
\]

und mit $S$ die im \S.\,79 betrachtete tern\"are Congruenzgruppe,
und setzen

\[
R = S\,Q,
\]
%CONT>>
